\clearpage\thispagestyle{empty}
\pdfbookmark[0]{Quellen zu Kometenbericht (September 1827)}{comet1827-sources}
\begin{center}\textbf{Quellen zu Kometenbericht (September 1827)}\end{center}
\noindent Diese KI-Übertragung ins heutige Deutsch folgt \textit{Comet} in \href{https://gdz.sub.uni-goettingen.de/id/PPN236007467}{\textit{Werke}, Band~VI} (Göttingen, 1874), S.~459--460. Heinrich Christian Schumacher berichtet über Gauß’ Meridianbeobachtungen und die von Peters berechnete Bahn. Der Erstdruck erschien im September 1827 in den \textit{Astronomischen Nachrichten}, Band~VI, Nr.~123, \href{https://api.digitale-sammlungen.de/iiif/image/v2/bsb10049431_00038/full/1550,/0/default.jpg}{Sp.~43--44}. Der gebundene Band der Bayerischen Staatsbibliothek trägt die Jahreszahl 1828.
\par\smallskip\noindent Die Werkausgabe gibt für die letzte Zeitangabe 47~27.0 und druckt nach der 5 in der Bahnneigung ein nach rechts offenes Zeichen, das hier als aufrechtes \textup{c} wiedergegeben wird. Im Erstdruck steht an diesen Stellen 47~27,7 beziehungsweise 50. Er liest außerdem \textit{Perihels} statt \textit{Prihels} und zeichnet mit \textit{S.} Die Übertragung folgt den Zahlen, dem Zeichen des aufsteigenden Knotens, der Tabellenanordnung und der ausgeschriebenen Unterschrift der Werkausgabe. KI-Übertragung und KI-Satz.
\clearpage
